%%%%%%%%%%     PREÁMBULO      %%%%%%%%%%%%%%%

%%%%%%%%%   Non modificar este preámbulo, salvo que se sepa ben o que se está a facer   %%%%%%%%%

%%%%%%  Traballando nun Mac, se usas TeXShop, en  Preferencias -> Código -> codificación de  caracteres (UTF-8)  %%%%

\documentclass[11pt,a4paper]{book}

% Formato do documento
\usepackage[papersize={210mm,297mm},lmargin=2.5cm,rmargin=2.5cm,top=3.5cm,bottom=3.5cm]{geometry}   % Marxes
\renewcommand{\baselinestretch}{1.25}                                                               % Fixamos o interliñado
\setlength{\parskip}{8pt}                                                                           % Separación entre parágrafos


\usepackage{amsmath,amsfonts,amssymb,amsthm}
\usepackage{latexsym,pdfsync,xcolor,graphicx}
\usepackage[nottoc]{tocbibind}
\usepackage[T1]{fontenc}
\usepackage[utf8]{inputenc}					
\usepackage[spanish]{babel}			
\usepackage{tikz}
\usepackage{appendix}
\usepackage{hyperref}








% Cabeceira
\usepackage{fancyhdr}
\pagestyle{fancy}
\fancyhf{}
\renewcommand{\chaptermark}[1]{\markboth{\textbf{\thechapter. #1}}{}} % Formato para o capítulo: N. Nome
\renewcommand{\sectionmark}[1]{\markright{\textbf{\thesection. #1}}}  % Formato para a sección: N.M. Nome

\fancyhead[LO]{\rightmark}           % Nas páxinas impares, parte esquerda do encabezado, aparecerá o nome do capítulo
\fancyhead[RE]{\leftmark}            % Nas páxinas pares, parte dereita do encabezado, aparecerá o nombre da sección
\fancyhead[RO,LE]{\textbf{\thepage}} % Números de páxina nas esquinas dos encabezados




%%%   \pagestyle{plain}   Cando os títulos dos capítulos son moi longos, a opción por defecto das cabeceiras pode ser inapropiada;
%%%  nestes e noutros casos semellantes, pódese usar a opción "plain", que suprime as cabeceiras, deixando soamente a numeración das páxinas.


%%%%%%%      NON ESCRIBIR AQUÍ  (COMANDO PORTADA) %%%%%%%%%%%

\newcommand{\portada}[3]{
\thispagestyle{empty}

\definecolor{azulUSC}{RGB}{13,38,120}

\includegraphics[width=13cm]{logomath.pdf}

\vspace*{2cm}

\centerline{{\Large\bf  Traballo Fin de Grao}}

\vspace{2.6cm}

{\center{\Huge\bfseries\color{azulUSC} #1}\par}		

\vspace{1cm}

{\Large

\centerline{#2}		

\vspace{6.4cm}

\centerline{\sf #3}        

\vspace{1.25cm}

\centerline{\sf UNIVERSIDADE DE SANTIAGO DE COMPOSTELA}
}

\enlargethispage{1cm}
\clearpage
}

%%%%%%%      NON ESCRIBIR AQUÍ  (COMANDO PRIMEIRA PÁXINA) %%%%%%%%%%%

\newcommand{\primeira}[3]{
\thispagestyle{empty}
\enlargethispage{1cm}
\begin{center}\Large

{\sf GRAO DE MATEM\'ATICAS}

\bigskip
{\bfseries Traballo Fin de Grao}
\vspace{4cm}

{\bfseries\Huge #1} 		
\vspace{1.2cm}

{#2}		\par		     
\end{center}

\vspace{6.5cm}

\begin{center}
\Large
{#3}	\par			    

\vspace{1cm}
{\sf UNIVERSIDADE DE SANTIAGO DE COMPOSTELA}
\end{center}
\clearpage
}

%%%%%%%    Teoremas   %%%%%%%%%

\newtheorem{theorem}{Teorema}[chapter]
\newtheorem{proposition}[theorem]{Proposición}
\newtheorem{lemma}[theorem]{Lema}
\newtheorem{corollary}[theorem]{Corolario}
\newtheorem{algorithm}[theorem]{Algoritmo}
\newtheorem{axiom}[theorem]{Axioma}
\newtheorem{case}[theorem]{Caso}
\newtheorem{conclusion}[theorem]{Conclusión}
\newtheorem{condition}[theorem]{Condición}
\newtheorem{conjetura}[theorem]{Conjetura}				
\newtheorem{criterion}[theorem]{Criterio}
\newtheorem{ejercicio}[theorem]{Ejercicio}	
			
\theoremstyle{definition}
\newtheorem{definition}[theorem]{Definición}
\newtheorem{ejemplo}[theorem]{Ejemplo}				
\newtheorem{agradecimientos}[theorem]{Agradecimientos}	

\theoremstyle{remark}
\newtheorem{remark}[theorem]{Observación}
\newtheorem{notation}[theorem]{Notación}
\newtheorem{problem}[theorem]{Problema}
\newtheorem{solution}[theorem]{Solución}

\addto\captionsspanish{
	\renewcommand{\contentsname}{\'Indice} 
	\renewcommand\appendixname{Anexo}
	\renewcommand\appendixpagename{Anexos}
}
%%%%%%%%%%%%%%%%%%   Fin do Preámbulo   %%%%%%%%%%%%%%%%




\begin{document}
\frontmatter

%%%%%%%%%%%      PORTADA      %%%%%%%%%%%%

\portada
{Resolución numérica del problema no lineal de mínimos cuadrados.\vspace{5pt}
Aplicaciones a la estimación de parámetros de modelos matemáticos.}		%  escribir aquí o  Título
{Dídac Blanco Morros}		%  escribir aquí o/a  Autor/a
{Curso Acad\'emico}		  	%  escribir aquí a Data de presentación, na forma:  mes, ano

%%%%%%%%     PRIMEIRA PÁXINA    %%%%%%%%%%%%%%%%%%

\mbox{}
\thispagestyle{empty}
\clearpage

\setcounter{page}{1}

\primeira
{Resolución numérica del problema no lineal de mínimos cuadrados. 
Aplicaciones a la estimación de parámetros de modelos matemáticos.}		%  escribir aquí o  Título
{Dídac Blanco Morros}		%  escribir aquí o/a  Autor/a
{Febrero, 2022}		  	%  escribir aquí a Data de presentación, na forma:  mes, ano

\thispagestyle{empty}

%%%%%%%%%%%      PROPOSTA DE TRABALLO      %%%%%%%%%%%%%%%%

\vspace*{.5 cm}

\chapter*{Trabajo propuesto}

\vspace{1 cm}

\begingroup
\renewcommand*{\arraystretch}{2.2}
\begin{tabular}{|l|}
	\hline
	
	{\bf \'Area de Co\~necemento:  \/ }\\ \hline
	\begin{minipage}{11.5cm}
		{\vspace*{.2cm}
			\bf T\'{\i}tulo:   \/ \bf
			\vspace{.2cm}}
	\end{minipage}\\ \hline
	\bf Breve descrici\'on do contido\\ \hline
	\begin{minipage}{11.5cm}
		{\vspace*{.2cm}

		
		\vspace{.2cm}}
	\end{minipage}\\ \hline
	{\bf Recomendaci\'ons}\\ \hline
	\begin{minipage}{11.5cm}
		{\vspace*{.2cm}


		\vspace{.2cm}}
	\end{minipage}\\ \hline
	{\bf Outras observaci\'ons}\\ \hline
	\begin{minipage}{11.5cm}
		{\vspace*{.2cm}


			\vspace{.2cm}}
	\end{minipage}\\ \hline
\end{tabular}
\endgroup

\clearpage

\thispagestyle{empty}

\tableofcontents

\clearpage

\thispagestyle{empty}

\mbox{}

\clearpage

%%%%%%%%%     RESUMO---RESUMEN---ABSTRACT     %%%%%%%%%%%%%

\phantomsection
\addcontentsline{toc}{chapter}{Resumen}		
\chapter*{}

\section*{Resumen}

% escribir aquí resumo en español




\vspace{1.5cm}

\section*{Abstract}



% escribir aquí resumo en inglés



\clearpage

\thispagestyle{empty}

%%%%%%%%%    INTRODUCIÓN     %%%%%%%%%%%%%%%%%%
%%%%%%%%%    De non incluír unha INTRODUCIÓN, suprímese    %%%%%%%%%%

\phantomsection
\chapter*{Introducci\'on}\addcontentsline{toc}{chapter}{Introducci\'on}

\markboth{INTRODUCCI\'ON}{INTRODUCCI\'ON}

%%%%%%%%%%%     fin da INTRODUCIÓN     %%%%%%%%%%%%%

\mainmatter

%%%%%%%%%%%%%%%%%%%%%%%%%%%%%%%%%%%%%%%%%%%%%%%%%%%%%				S
%%%%%%%%%%%%%%%%%%%%%%%%%%%%%%%%%%%%%%%%%%%%%%%%%%%%%				 T
%%%%%%%%%%%%%%%%%%%%%%%%%%%%%%%%%%%%%%%%%%%%%%%%%%%%%				  A
%%%%%%%%%%%%%%%%%%%%%%%%%%%%%%%%%%%%%%%%%%%%%%%%%%%%%				   R
%%%%%%%%     Póñense os Capítulos    %%%%%%%%%%%%%%%				    T

\chapter{Título del Capítulo 1}



\section{Fundamentos de la optimización sin restricciones}

El problema de los mínimos cuadrados es un caso particular de optimización sin restricciones, y es por ello que comenzaremos introduciendo sus fundamentos. Ya que el problema de mínimos cuadrados es usado en multitud de campos para estimar parámetros, este es de los más utilizados dentro de los problemas de optimización sin restricciones.

Un problema de optimización sin restricciones tiene la forma 
\begin{equation}
	\min_{x}f\left(x\right)
	\label{eq:minf}
\end{equation}
donde $x\in\mathbb{R}^{n}$ y $f : \mathbb{R}^{n} \rightarrow \mathbb{R}$ es continuamente diferenciable, la llamamos \textbf{función objetivo}. La dificultad de un problema como este viene de no conocer el comportamiento global de $f$, normalmente solo disponemos de la evaluación de f en algunos puntos, y a lo mejor de algunas de sus derivadas. El trabajo de los algoritmos de optimización es identificar la solución sin usar demasiado tiempo ni almacenamiento computacional.

Notar que podemos usar la formulación (\ref{eq:minf}) para referirnos a tanto a los problemas de minimización como de maximización, basta sustituir $f$ por $-f$. 

% W. Sun es más formal en este caso para definir ambos mínimos 
Tenemos dos tipos de mínimos.
Un punto $x^{*}$ se dice \textbf{mínimo global} si $f(x^{*})\leq f(x)$ para todo $x\in\mathbb{R}^{n}$. Como no se suele tener un conocimiento a gran escala de $f$ debido a su coste, la mayoría de algoritmos solo encuentran mínimos locales, lo cual es suficiente para muchos casos prácticos. Un punto $x^{*}$ se dice \textbf{mínimo local} si existe una vecinidad $\mathcal{V}$ de $x^{*}$ tal que $f(x^{*})\leq f(x)$ para todo $x\in\mathcal{V}$.

Aún así, los algoritmos para encontrar mínimos globales se suelen construír a partir de una secuencia de otros algoritmos de optimización local. También podemos aprovechar características fáciles de detectar en la función objetivo, como la convexidad, que nos asegura que un mínimo local será también global.

%% REPASO DE ALGORITMOS

\section{Sobre los algoritmos}

%% histórico??
Todo algoritmo de optimización sin restricciones comienza con un punto de partida, denotado normalmente como $x_{0}$. Aunque generalmente el usuario introduce una estimación razonable, el punto puede ser elegido por el algoritmo, tanto de forma sistemática como aleatoria. El algoritmo itera sobre $x_{0}$, creando una sucesión $\left\{x_k\right\}_{k=0}^n$ la cual termina cuando se cumpla cierto criterio.



\chapter{Título del Capítulo 2}




\chapter{Título del Capítulo 3}




%%%%%%%%     Póñense a bibliografía, Apéndices e o índice alfabético     %%%%%%%%
\appendix
\renewcommand{\thechapter}{\Roman{chapter}}
\chapter{Título del Anexo I}


\chapter{Título del Anexo II}

\backmatter

\begin{thebibliography}{99}

%%%%%%%%----Exemplo de entradas bibliográficas:
%
\bibitem{HiSt} Hilton, P. J. y Stammbach, U. A. (1997). \emph{A course in homological algebra}, 2nd ed., Graduate Texts in Mathematics, 4, Springer-Verlag,  New York.
\bibitem{Lo} Loday, J.-L. (1996). \emph{Künneth-style formula for the homology of Leibniz algebras}, Math. Z., \textbf{221}, 41--47.
%
\end{thebibliography}




\end{document}


